(1) 取单段球面规范化测度 \(d\Omega/(4\pi)\)，记 \(x=\beta fb=fb/(k_BT)\)。单段配分函数
\[
z_1(x)=\frac12\int_{-1}^{1}e^{xu}du=\frac{\sinh x}{x},\quad
Z=z_1(x)^N,\quad
\Psi(f,T)=-Nk_BT\ln\frac{\sinh x}{x}.
\]
对广义力求导得
\[
z=-\left(\frac{\partial\Psi}{\partial f}\right)_T
=Nb\,\mathcal L(x),\quad \mathcal L(x)=\coth x-\frac1x.
\]
这里 \(\mathcal L\) 是 Langevin 函数，取 \(x=0\) 的连续极限为零。转动自由度以外的平动及链段振动配分函数不纳入构型模型。

(2) 直接微分有
\[
\left(\frac{\partial z}{\partial f}\right)_T
=\frac{Nb^2}{k_BT}\left(\frac1{x^2}-\frac1{\sinh^2x}\right)>0
\quad(x\ne0).
\]
严格正性来自 \(|\sinh x|>|x|\)，原点连续极限为 \(Nb^2/(3k_BT)\)。由
\(\mathcal L(x)=x/3-x^3/45+2x^5/945+\cdots\)，
可反演为 \(x=3y+\frac95y^3+O(y^5)\)。故
\[
f=\frac{k_BT}{b}\left(3y+\frac95y^3+O(y^5)\right).
\]
在定伸长系综，\(dA=f\,dz=Nk_BT\,x\,dy\)，从自然态 \(y=0\) 积分得
\[
A(y,T)-A(0,T)
=Nk_BT\left(\frac32y^2+\frac9{20}y^4+O(y^6)\right).
\]
这独立核验了高斯链弹簧常数 \(3k_BT/(Nb^2)\)。

(3) 在大 \(N\) 的热力学极限，定伸长势是 Legendre 变换
\[
A(z,T)=\Psi(f,T)+fz=Nk_BT
\left[x\mathcal L(x)-\ln\frac{\sinh x}{x}\right],
\quad y=\mathcal L(x).
\]
对 \(z\) 的导数为
\(\partial A/\partial z=f\)，且
\[
\frac{\partial^2A}{\partial z^2}
=\frac{k_BT}{Nb^2\,\mathcal L'(x)}>0.
\]
当 \(y\to1^-\) 时，\(\mathcal L(x)=1-1/x+O(e^{-2x})\)，
故 \(x\sim(1-y)^{-1}\)，\(f\sim k_BT/[b(1-y)]\)，并且
\[
\frac{A}{Nk_BT}
=\ln\frac2{1-y}-1+o(1).
\]
自由能的对数发散表达链段取向被压缩到零测度集合的熵损失，而非额外的谐弹性键能。

(4) 因为在固定 \(y\) 时 \(x=\mathcal L^{-1}(y)\) 不依赖温度，
\[
\Delta S(z,T)=-\left(\frac{\partial\Delta A}{\partial T}\right)_z
=-Nk_B\left[x y-\ln\frac{\sinh x}{x}\right],\qquad
\Delta U=\Delta A+T\Delta S=0,\quad C_z=0.
\]
因此温度上升时要保持同一伸长，拉力按 \(T\) 成正比增长，这是熵弹性；不能把力控哈密顿量中的外势 \(-fz\) 错认作链内部的弹性储能。\(f\to0\) 或 \(T\to\infty\) 的固定 \(f\) 极限均给 \(z\sim Nb^2 f/(3k_BT)\to0\)。严格界限 \(|z|<Nb\)；有限 \(N\) 的精确定伸长分布需做 \(N\) 段取向卷积，本题 Legendre 形式描述其主阶大偏差自由能，非声称有限 \(N\) 的逐态精确等式。
